% Hippocampus Hypergraph Migration: From Flat Memory to Interconnected Knowledge
% Samuel Frerichs, Michael Clark — Are-Self Research
\documentclass[man,12pt,floatsintext]{apa7}

\usepackage[utf8]{inputenc}
\usepackage[T1]{fontenc}
\usepackage{amsmath,amssymb,amsfonts}
\usepackage{graphicx}
\usepackage{booktabs}
\usepackage{hyperref}
\usepackage{algorithm}
\usepackage{algpseudocode}
\usepackage{listings}
\usepackage{xcolor}

\definecolor{codebg}{rgb}{0.95,0.95,0.97}
\definecolor{codegreen}{rgb}{0.2,0.6,0.2}
\definecolor{codegray}{rgb}{0.5,0.5,0.5}
\definecolor{codepurple}{rgb}{0.58,0,0.82}

\lstset{
  backgroundcolor=\color{codebg},
  basicstyle=\ttfamily\small,
  breaklines=true,
  captionpos=b,
  commentstyle=\color{codegreen},
  keywordstyle=\color{codepurple},
  numberstyle=\tiny\color{codegray},
  stringstyle=\color{codegreen},
  frame=single,
  numbers=left,
}

\usepackage[style=apa,backend=biber]{biblatex}
\addbibresource{../../templates/references.bib}

\title{Hippocampus Hypergraph Migration: From Flat Memory to Interconnected Knowledge}
\shorttitle{HIPPOCAMPUS HYPERGRAPHS}
\author{Samuel Frerichs\textsuperscript{1} and Michael Clark\textsuperscript{2}}
\affiliation{\textsuperscript{1}University of Pittsburgh at Applied Sciences \\
\textsuperscript{2}Independent Researcher, scipraxian}

\abstract{
Are-Self's Hippocampus module stores long-term memories as vector-embedded Engrams with
flat many-to-many relationships to sessions, identities, and tags. While effective for
single-fact retrieval and cosine-similarity deduplication, this flat structure cannot
represent the rich interconnections between memories that characterize human episodic and
semantic recall. We propose migrating the Hippocampus from flat M2M relations to a
hypergraph structure, where Engrams become nodes and typed hyperedges represent
relationships such as causation, temporal sequence, contradiction, elaboration, and
generalization. We formalize the hypergraph memory model, describe a migration path that
preserves backward compatibility with existing vector search, and analyze the expected
improvements in contextual memory retrieval, knowledge consolidation, and reasoning
session quality. This work extends the neuro-mimetic design principle: just as the
biological hippocampus encodes not just facts but their interconnections, our software
Hippocampus should do the same.
}

\keywords{knowledge graphs, hypergraph memory, vector embeddings, episodic memory,
semantic memory, neuro-mimetic architecture}

\begin{document}
\maketitle

\section{Introduction}

The human hippocampus does not store memories as isolated facts. It encodes episodes---
temporally ordered sequences of events with causal, spatial, and emotional
interconnections. Retrieval is associative: recalling one memory activates related
memories through these connections, enabling reasoning by analogy, contradiction
detection, and knowledge synthesis.

Are-Self's current Hippocampus module uses a flat model: each Engram is a named fact with
a 768-dimensional vector embedding, provenance metadata (sessions, spikes, identities),
and tags. Retrieval is pure vector similarity: given a query, return the $K$ most
similar engrams. This works well for factual recall but cannot answer questions like
``What happened after the deployment failed?'' or ``Does this contradict what we learned
yesterday?''

We propose extending the Hippocampus with a hypergraph layer that preserves vector search
while adding typed relationships between engrams.

\section{Related Work}

\subsection{Knowledge Graphs and Memory Systems}

MemGPT \parencite{} introduces hierarchical memory tiers for LLM agents. Mem0 provides
persistent memory with retrieval. These systems focus on storage and retrieval optimization
but do not model inter-memory relationships.

\subsection{Hypergraph Databases}

Unlike standard graphs where edges connect exactly two nodes, hyperedges connect arbitrary
sets of nodes. This is natural for memory relationships: ``these three facts together
caused this outcome'' requires a hyperedge, not pairwise edges.

\subsection{Biological Hippocampal Models}

Computational neuroscience models of hippocampal memory consolidation (complementary
learning systems theory) distinguish between rapid episodic encoding and gradual semantic
consolidation. Our hypergraph migration maps these processes to software operations.

\section{Current Architecture}

\subsection{The Engram Model}

% TODO: Describe current flat model in detail
% - Fields: name (hash), description (fact), vector (768-dim), relevance_score
% - M2M: sessions, source_turns, spikes, creator, identity_discs, tags, tasks
% - Deduplication: cosine similarity >= 0.90 rejects save
% - Auto-vectorization: nomic-embed-text via Ollama

\subsection{Retrieval Pipeline}

% TODO: Describe current vector-similarity retrieval
% - Query vector computed at retrieval time
% - Top-K nearest neighbors by cosine distance
% - Filtered by creator/identity_disc access
% - Injected into HISTORY phase of addon pipeline

\subsection{Limitations of Flat Memory}

% TODO: Concrete examples of limitations
% - Cannot answer "what happened next?"
% - Cannot detect contradictions between memories
% - Cannot identify memory clusters (related facts)
% - Cannot represent causal chains

\section{Hypergraph Memory Model}

\subsection{Formal Definition}

Let $\mathcal{H} = (V, E)$ be a hypergraph where $V$ is the set of Engrams and $E$ is a
set of typed hyperedges. Each hyperedge $e \in E$ has:

\begin{itemize}
  \item A type $\tau(e) \in \{$TEMPORAL, CAUSAL, CONTRADICTS, ELABORATES,
        GENERALIZES, CLUSTERS$\}$
  \item A source set $S(e) \subseteq V$
  \item A target set $T(e) \subseteq V$
  \item A confidence score $c(e) \in [0, 1]$
  \item A creation context (which session/turn detected the relationship)
\end{itemize}

\subsection{Edge Type Semantics}

% TODO: Define each edge type formally
% - TEMPORAL: S happened before T
% - CAUSAL: S caused T
% - CONTRADICTS: S and T assert incompatible facts
% - ELABORATES: T provides detail on S
% - GENERALIZES: T is a generalization of S
% - CLUSTERS: S and T belong to a semantic cluster

\subsection{Automatic Edge Detection}

% TODO: Algorithm for detecting relationships during engram save
% - Temporal: session timestamps
% - Causal: LLM-assessed from context
% - Contradicts: high vector similarity + semantic negation
% - Elaborates: vector similarity + substring/expansion detection
% - Generalizes: subsumption checking
% - Clusters: community detection on similarity graph

\section{Migration Strategy}

\subsection{Backward Compatibility}

% TODO: How to migrate without breaking existing vector search
% - Additive schema: new HyperedgeModel alongside existing Engram
% - Vector search remains primary retrieval path
% - Hypergraph traversal as optional enrichment step
% - Gradual edge population via background task

\subsection{Database Considerations}

% TODO: pgvector + hypergraph edge table in PostgreSQL
% - vs. dedicated graph database (Neo4j)
% - Performance implications
% - Django ORM integration

\section{Enhanced Retrieval}

\subsection{Context-Enriched Recall}

Given a query, the enhanced retrieval pipeline:

\begin{enumerate}
  \item Standard vector-similarity top-$K$ retrieval
  \item Hypergraph traversal from retrieved nodes (1-2 hops)
  \item Edge-type filtering based on query intent
  \item Re-ranking by combined vector similarity + graph centrality
  \item Return enriched memory context with relationship annotations
\end{enumerate}

\subsection{Contradiction Detection}

% TODO: Algorithm for detecting contradictions between new and existing memories

\subsection{Knowledge Consolidation}

% TODO: Periodic background process that:
% - Identifies engram clusters
% - Generates summary engrams (GENERALIZES edges)
% - Prunes redundant edges
% - Mirrors biological memory consolidation during "sleep"

\section{Evaluation Plan}

% TODO: Experimental design
% - Retrieval relevance: flat vs. hypergraph-enriched
% - Contradiction detection accuracy
% - Reasoning session quality (turns to solution, memory utilization)
% - Computational overhead of hypergraph operations

\section{Discussion}

The hypergraph migration extends Are-Self's neuro-mimetic principle into the memory domain.
Just as the biological hippocampus maintains rich interconnections between memories that
enable associative recall, our software Hippocampus will support relationship-aware
retrieval that contextualizes individual facts within the broader knowledge landscape.

\subsection{Future Work}

Sleep-mode consolidation that runs during idle periods to strengthen frequently-accessed
edges, merge redundant clusters, and generate abstract summaries. Multi-agent memory
sharing through federated hypergraph segments. Emotional tagging via a software
Amygdala module.

\section{Conclusion}

We have proposed a concrete migration path from Are-Self's flat Engram model to a
hypergraph-based memory system that preserves vector search while adding typed
relationships between memories. The hypergraph structure enables contextual recall,
contradiction detection, and knowledge consolidation---capabilities that the flat model
cannot support. This migration is planned for implementation in collaboration between the
University of Pittsburgh at Applied Sciences and the Are-Self open-source project.

\printbibliography

\end{document}
